\documentclass[10pt,a4paper]{amsart}
\usepackage{amsmath,amssymb,amsthm,graphicx}
% For commutative diagrams we use Xy-pic
\usepackage[all]{xy}

% We use 2cell for 2-commutative diagrams.
\xyoption{2cell}
% Theorem environments.
%
\theoremstyle{plain}
\newtheorem{theorem}[subsection]{Theorem}
\newtheorem{proposition}[subsection]{Proposition}
\newtheorem{lemma}[subsection]{Lemma}

\theoremstyle{definition}
\newtheorem{definition}[subsection]{Definition}
\newtheorem{example}[subsection]{Example}
\newtheorem{exercise}[subsection]{Exercise}
\newtheorem{situation}[subsection]{Situation}

\theoremstyle{remark}
\newtheorem{remark}[subsection]{Remark}
\newtheorem{remarks}[subsection]{Remarks}

\numberwithin{equation}{subsection}

% Macros
%
\def\lim{\mathop{\mathrm{lim}}\nolimits}
\def\colim{\mathop{\mathrm{colim}}\nolimits}
\def\Spec{\mathop{\mathrm{Spec}}}
\def\Hom{\mathop{\mathrm{Hom}}\nolimits}
\def\Ext{\mathop{\mathrm{Ext}}\nolimits}
\def\SheafHom{\mathop{\mathcal{H}\!\mathit{om}}\nolimits}
\def\SheafExt{\mathop{\mathcal{E}\!\mathit{xt}}\nolimits}
\def\Sch{\mathit{Sch}}
\def\Mor{\mathop{\mathrm{Mor}}\nolimits}
\def\Ob{\mathop{\mathrm{Ob}}\nolimits}
\def\Sh{\mathop{\mathit{Sh}}\nolimits}
\def\NL{\mathop{N\!L}\nolimits}
\def\CH{\mathop{\mathrm{CH}}\nolimits}
\def\proetale{{pro\text{-}\acute{e}tale}}
\def\etale{{\acute{e}tale}}
\def\QCoh{\mathit{QCoh}}
\def\Ker{\mathop{\mathrm{Ker}}}
\def\Im{\mathop{\mathrm{Im}}}
\def\Coker{\mathop{\mathrm{Coker}}}
\def\Coim{\mathop{\mathrm{Coim}}}

% Boxtimes
%
\DeclareMathSymbol{\boxtimes}{\mathbin}{AMSa}{"02}

%
% Macros for moduli stacks/spaces
%
\def\QCohstack{\mathcal{QC}\!\mathit{oh}}
\def\Cohstack{\mathcal{C}\!\mathit{oh}}
\def\Spacesstack{\mathcal{S}\!\mathit{paces}}
\def\Quotfunctor{\mathrm{Quot}}
\def\Hilbfunctor{\mathrm{Hilb}}
\def\Curvesstack{\mathcal{C}\!\mathit{urves}}
\def\Polarizedstack{\mathcal{P}\!\mathit{olarized}}
\def\Complexesstack{\mathcal{C}\!\mathit{omplexes}}
% \Pic is the operator that assigns to X its picard group, usage \Pic(X)
% \Picardstack_{X/B} denotes the Picard stack of X over B
% \Picardfunctor_{X/B} denotes the Picard functor of X over B
\def\Pic{\mathop{\mathrm{Pic}}\nolimits}
\def\Picardstack{\mathcal{P}\!\mathit{ic}}
\def\Picardfunctor{\mathrm{Pic}}
\def\Deformationcategory{\mathcal{D}\!\mathit{ef}}
\usepackage{cleveref}
\crefname{originallemma}{original lemma}{original lemmas}
\Crefname{originallemma}{Original lemma}{Original lemmas}
\makeatletter
\newlabel{algebra-lemma-dominate-by-dimension-1}{{00P8}{source}{}{sourcepointer-1}{}}
\newlabel{algebra-lemma-dominate-by-dimension-1@cref}{{[originallemma][0][]00P8}{[1][0][]source}{}{}{}}
\makeatother
\crefname{originallemma}{original lemma}{original lemmas}
\Crefname{originallemma}{Original lemma}{Original lemmas}
\makeatletter
\newlabel{algebra-lemma-integral-overring-surjective}{{00GQ}{source}{}{sourcepointer-2}{}}
\newlabel{algebra-lemma-integral-overring-surjective@cref}{{[originallemma][0][]00GQ}{[1][0][]source}{}{}{}}
\makeatother
\crefname{originallemma}{original lemma}{original lemmas}
\Crefname{originallemma}{Original lemma}{Original lemmas}
\makeatletter
\newlabel{algebra-lemma-surjective-spec-radical-ideal}{{00FI}{source}{}{sourcepointer-3}{}}
\newlabel{algebra-lemma-surjective-spec-radical-ideal@cref}{{[originallemma][0][]00FI}{[1][0][]source}{}{}{}}
\makeatother
\title{A nonzerodivisor in a one-dimensional Noetherian local ring is only boundedly divisible by a parameter in finite extensions}
\author{The Stacks Project Authors}
\date{stacks-project:chow.tex:local-snapshot-20261011; GNU FDL 1.2 or later}
\begin{document}
\maketitle
\noindent\textit{Editorial note.} All delivered mathematical text is the author's own, copied from the pinned Stacks Project chapter file chow.tex; no statement or proof step was rewritten, reordered or removed.\par\medskip
\noindent\textit{Editorial note.} The preamble is the chapter's own native preamble excerpt used for this source: the Xy-pic setup at lines 14-18 and the theorem declarations with the mathematical macros at lines 154-217 of sources/stacks-chow/preamble.tex. The quoted class line and the xr-hyper externaldocument block are not included because the assembler supplies a plain amsart wrapper; no macro was invented or redefined.\par\medskip
\noindent\textit{Editorial note.} All three referenced results live in earlier chapters, so each locator points at the verbatim line of the project's own tag list staged in the snapshot (sources/stacks-chow/tags.txt) that carries the canonical tag.\par\medskip
\noindent\textit{Allowed standard prerequisites (editorial).}\begin{itemize}
\item Dominating a one-dimensional Noetherian domain by a discrete valuation ring. Algebra, Lemma labelled algebra-lemma-dominate-by-dimension-1 (canonical Stacks tag 00P8); invoked at proof lines 615-617 and again at proof lines 633-634. Supplies the discrete valuation ring A dominating R/q with the same fraction field, and later the ring A''' dominating A''\_m''; the valuation estimate is the counted content.
\item Surjectivity of spectra for integral and for radical extensions. Algebra, Lemma algebra-lemma-integral-overring-surjective (00GQ, proof lines 625-626) and Lemma algebra-lemma-surjective-spec-radical-ideal (00FI, proof lines 627-628). Supply that Spec(R') -> Spec(R) and Spec(A') -> Spec(A) are surjective, which is what produces the prime q' of A' lying over (0) of A.
\end{itemize}\medskip
\ifdefined\hypertarget\hypertarget{sourcepointer-1}{}\fi
\noindent\textit{Source reference.} Original lemma 00P8: sources/stacks-chow/tags.txt, source lines 922--922. Cross-chapter standard prerequisite invoked at proof lines 615-617 and again at 633-634 to obtain the two discrete valuation rings; the chain of base changes and the valuation comparison are retained in the counted span.\par\medskip
\ifdefined\hypertarget\hypertarget{sourcepointer-2}{}\fi
\noindent\textit{Source reference.} Original lemma 00GQ: sources/stacks-chow/tags.txt, source lines 616--616. Cross-chapter standard prerequisite invoked once at proof lines 625-626 for surjectivity of Spec(R') -> Spec(R).\par\medskip
\ifdefined\hypertarget\hypertarget{sourcepointer-3}{}\fi
\noindent\textit{Source reference.} Original lemma 00FI: sources/stacks-chow/tags.txt, source lines 572--572. Cross-chapter standard prerequisite invoked once at proof lines 627-628 for surjectivity of Spec(A') -> Spec(A); the choice of the prime q' and the rest of the argument are retained in the counted span.\par\medskip
\begin{lemma}
\label{lemma-not-infinitely-divisible}
Let $(R, \mathfrak m)$ be a Noetherian local ring of dimension $1$.
Let $a, b \in R$ be nonzerodivisors with $a \in \mathfrak m$.
There exists an integer $n = n(R, a, b)$ such that for a finite ring
extension $R \subset R'$ if $b = a^m c$ for some $c \in R'$, then $m \leq n$.
\end{lemma}

\begin{proof}
Choose a minimal prime $\mathfrak q \subset R$. Observe that
$\dim(R/\mathfrak q) = 1$, in particular $R/\mathfrak q$ is not a field.
We can choose a discrete valuation ring $A$ dominating $R/\mathfrak q$
with the same fraction field, see
Algebra, Lemma \ref{algebra-lemma-dominate-by-dimension-1}. Observe that
$a$ and $b$ map to nonzero elements of $A$ as nonzerodivisors in $R$
are not contained in $\mathfrak q$. Let $v$ be the discrete valuation on $A$.
Then $v(a) > 0$ as $a \in \mathfrak m$.
We claim $n = v(b)/v(a)$ works.

\medskip\noindent
Let $R \subset R'$ be given. Set $A' = A \otimes_R R'$.
Since $\Spec(R') \to \Spec(R)$ is surjective
(Algebra, Lemma \ref{algebra-lemma-integral-overring-surjective})
also $\Spec(A') \to \Spec(A)$ is surjective
(Algebra, Lemma \ref{algebra-lemma-surjective-spec-radical-ideal}).
Pick a prime $\mathfrak q' \subset A'$ lying over $(0) \subset A$.
Then $A \subset A'' = A'/\mathfrak q'$ is a finite extension of rings
(again inducing a surjection on spectra).
Pick a maximal ideal $\mathfrak m'' \subset A''$
lying over the maximal ideal of $A$ and a discrete valuation ring
$A'''$ dominating $A''_{\mathfrak m''}$ (see lemma cited above).
Then $A \to A'''$ is an extension of discrete valuation rings
and we have $b = a^m c$ in $A'''$. Thus $v'''(b) \geq mv'''(a)$.
Since $v''' = ev$ where $e$ is the ramification index
of $A'''/A$, we find that $m \leq n$ as desired.
\end{proof}

\end{document}
